% Options for packages loaded elsewhere
\PassOptionsToPackage{unicode}{hyperref}
\PassOptionsToPackage{hyphens}{url}
%
\documentclass[
  10pt,
]{article}
\usepackage{amsmath,amssymb}
\usepackage{iftex}
\ifPDFTeX
  \usepackage[T1]{fontenc}
  \usepackage[utf8]{inputenc}
  \usepackage{textcomp} % provide euro and other symbols
\else % if luatex or xetex
  \usepackage{unicode-math} % this also loads fontspec
  \defaultfontfeatures{Scale=MatchLowercase}
  \defaultfontfeatures[\rmfamily]{Ligatures=TeX,Scale=1}
\fi
\usepackage{lmodern}
\ifPDFTeX\else
  % xetex/luatex font selection
\fi
% Use upquote if available, for straight quotes in verbatim environments
\IfFileExists{upquote.sty}{\usepackage{upquote}}{}
\IfFileExists{microtype.sty}{% use microtype if available
  \usepackage[]{microtype}
  \UseMicrotypeSet[protrusion]{basicmath} % disable protrusion for tt fonts
}{}
\makeatletter
\@ifundefined{KOMAClassName}{% if non-KOMA class
  \IfFileExists{parskip.sty}{%
    \usepackage{parskip}
  }{% else
    \setlength{\parindent}{0pt}
    \setlength{\parskip}{6pt plus 2pt minus 1pt}}
}{% if KOMA class
  \KOMAoptions{parskip=half}}
\makeatother
\usepackage{xcolor}
\usepackage[margin=0.85in]{geometry}
\setlength{\emergencystretch}{3em} % prevent overfull lines
\providecommand{\tightlist}{%
  \setlength{\itemsep}{0pt}\setlength{\parskip}{0pt}}
\setcounter{secnumdepth}{-\maxdimen} % remove section numbering
\usepackage{amsmath,amssymb,amsthm}
\usepackage{longtable}
\setlength{\emergencystretch}{3em}
\ifLuaTeX
  \usepackage{selnolig}  % disable illegal ligatures
\fi
\usepackage{bookmark}
\IfFileExists{xurl.sty}{\usepackage{xurl}}{} % add URL line breaks if available
\urlstyle{same}
\hypersetup{
  pdftitle={P16.2: retained polar data on proper nonprojective varieties},
  pdfauthor={Research draft},
  hidelinks,
  pdfcreator={LaTeX via pandoc}}

\title{P16.2: retained polar data on proper nonprojective varieties}
\usepackage{etoolbox}
\makeatletter
\providecommand{\subtitle}[1]{% add subtitle to \maketitle
  \apptocmd{\@title}{\par {\large #1 \par}}{}{}
}
\makeatother
\subtitle{Proof reconstruction, original sections, and a remaining
local-operator bridge}
\author{Research draft}
\date{5 October 2026}

\begin{document}
\maketitle

\section{Abstract and status}\label{abstract-and-status}

This paper gives a partial reconstruction of original P16.2. The ground
field is any algebraically closed trivially valued field of
characteristic zero, with no cardinality assumption. The original
carrier is any smooth integral proper variety of positive dimension. We
prove existence of an actual big line bundle on every such carrier and
of a retained polar presentation for every specified actual big rational
line bundle. We construct actual original metrics, prove their section
approximation and all-point descent, and repair the fixed-divisor cone
argument without the false general extended-valued flattening lemma. We
then construct a positive raw measure-valued pairing before choosing any
target, prove its mass and finite-energy transports, and prove the
corresponding conditional variational equation theory.

The completed local argument has a specific scope: decreasing finite
section maxima in a fixed actual ample bundle on the retained
modification, divided on original frame neighborhoods. It covers bounded
discontinuous endpoints and passes frame-unit identities through those
limits. What remains unproved here is the agreement of this construction
with the local contact/collar intersection product in the complete
approximation category required for an intrinsic original-carrier local
product. This is a compatibility and independence statement for the
actual weights constructed below, not a request for a theorem about
arbitrary unrelated local psh functions. Consequently this paper does
\textbf{not} claim a complete solution of P16.2. Its status is
\texttt{candidate=false}.

All results below are mathematical proof drafts. The supplied
actual-ample theorem G0 is the sole task grant, in exactly its actual
ample domain. Published background theorems are recorded separately.
Supplied historical reports and continuations supply no additional
grants. Executions, hashes, and PDF compilation have no mathematical
force.

\section{1. Exact domain, conventions, and the ample
grant}\label{exact-domain-conventions-and-the-ample-grant}

Fix an algebraically closed field \(k\) of characteristic zero with its
trivial absolute value. Let \(X\) be a smooth integral proper
\(k\)-variety of dimension \(n>0\). Projectivity of \(X\) is not
assumed. By an actual rational line bundle we mean an element of
\(\operatorname{Pic}(X)\otimes\mathbb Q\), with all section calculations
made after taking specified integral multiples. A big actual rational
bundle \(L\) means that an integral multiple \(F=qL\) satisfies \[
\limsup_{m\longrightarrow\infty}m^{-n}h^0(X,mF)>0.
\] No section ring is assumed finitely generated. The construction
retains its choices; it does not identify theories from different
retained presentations.

The replacement datum will be \[
\mathcal M=(L,\pi:Y\longrightarrow X,A,D,\iota),\qquad
\iota:\pi^*L\simeq A+D,                                      \tag{1.1}
\] where \(Y\) is smooth integral projective, \(\pi\) is a proper
birational morphism, \(A\) is an actual ample rational line bundle, and
\(D\) is an effective rational Cartier divisor. The line-bundle
isomorphism is part of the datum. All choices are made before a target
measure is considered. Put \[V=(A^n)>0.\] This is the raw mass of the
constructed pairing. We do not identify it with an unrestricted
minimal-singularity volume of \(L\), or with the possibly different
intersection number \((L^n)\).

\textbf{G0, used only on \((Y,A)\).} The task grants the genuine
finite-energy theory on geometrically unibranch integral projective
carriers with an actual ample rational bundle over the present fields.
In particular it grants the following on the smooth projective \(Y\):

\begin{itemize}
\tightlist
\item
  genuine \(A\)-psh potentials, finite fixed-bundle section
  approximation, divisorial determination, normalized compactness, and
  envelopes;
\item
  the genuine raw symmetric mixed intersection products for
  \(\mathcal E^1(Y,A)\), of mass \(V\), with locality, comparison,
  domination, and unbounded decreasing-net continuity;
\item
  integrability of finite-energy tests for mixed finite-energy products,
  and integrable weighted decreasing limits;
\item
  raw energy \(E_A\), its usc and monotone continuity properties, and
  its two-sided outer derivative in every direction in
  \(C^0(Y^{\rm an})\);
\item
  the intrinsic ample target condition, its equivalence to integrability
  of all actual \(\mathcal E^1(Y,A)\) potentials, actual variational
  attainment, and uniqueness modulo constants.
\end{itemize}

For a bounded genuine ample potential the granted raw energy has the
explicit formula \[E_A(\psi)=\frac1{n+1}\sum_{j=0}^n
\int_Y\psi\,(\omega_A+dd^c\psi)^j\wedge\omega_A^{n-j},\qquad
\omega_A=c_1(A)\text{ with its canonical zero metric}.\] For an upper
bounded unbounded genuine \(\psi\), it is the decreasing limit of these
energies for \(\max(\psi,-j)\); value \(-\infty\) is permitted. The
genuine finite-energy space consists exactly of potentials with finite
raw energy. These are the actual-ample definitions imported from G0,
before any transport to the original carrier. We use raw energy,
normalized by \[E_A(0)=0,\qquad E_A(\psi+c)=E_A(\psi)+Vc.\] The
probability calibration is \(e_A=V^{-1}E_A\). G0 is not a grant of
nonample varying-class closure, divisor division, nonprojective descent,
or a new local metric construction. Whenever its actual-ample
consequences are used below, the carrier is \(Y\) and the bundle is
\(A\), or another explicitly chosen actual ample bundle on \(Y\).

All analytic spaces are Berkovich analytifications. A Radon measure is a
regular Borel measure. Nets, compactness, and continuous test functions
are used without metrizability assumptions.

\section{2. Two approximation diagnostics and their
repairs}\label{two-approximation-diagnostics-and-their-repairs}

\subsection{2.1. The supplied counterexample is a proof, not an equation
obstruction}\label{the-supplied-counterexample-is-a-proof-not-an-equation-obstruction}

Let \(K\) be the disjoint union of an isolated point \(q\) and the
compact convergent sequence \(\{p_j:j\geq1\}\cup\{p_\infty\}\). Set
\(S=\mathbb N_{>0}\), \(T=\mathbb R\), and \(\rho(a)=1/a\). For
\(a,m\geq1\) and \(c\in\mathbb R\) let \[
f_{a,m,c}(q)=c,\quad f_{a,m,c}(p_a)=-\infty,\quad
f_{a,m,c}(x)=c-m\quad\text{at all remaining points}.
\] Each function is continuous into \(\mathbb R\cup\{-\infty\}\),
including at \(p_\infty\), and each family \(F_a\) is stable under
adding constants. For fixed \(a\), \(f_{a,m,0}\) decreases to the same
function \(g\): its value is zero at \(q\) and \(-\infty\) at every
other point. Thus \(g\in\widetilde F_{1/a}\) for every \(a\), and the
constant outer family \(g_a=g\) has parameters tending to zero.

Nevertheless \(\widetilde F_0\) is empty. In any decreasing net of
members of the \(F_a\), fix one term from \(F_a\). Every later term must
also be from \(F_a\), since its value at \(p_a\) must remain
\(-\infty\), whereas every member of every other family is finite there.
On that entire tail its parameter is \(1/a\), so the parameters cannot
converge to zero.

This disproves BJ v4 Lemma 4.6 in its printed generality. Its
strict-shift step fails at a common \(-\infty\) value; the corresponding
exponential difference is zero rather than strictly negative. The proof
does not disprove any specialized algebraic FS closure theorem, and
gives no counterexample to the desired equation theorem. We do not use
the lemma, or the nonample conclusions of BJ 4.15, 6.17, and 6.24
obtained through it.

\subsection{2.2. Safe finite fixed-bundle
flattening}\label{safe-finite-fixed-bundle-flattening}

Let \(B\) be an actual ample rational bundle on an integral projective
variety in the present field domain. Write \(F_B\) for finite,
real-valued, continuous actual FS section maxima \[
m^{-1}\max_j(\log\|s_j\|_0+\lambda_j),
\] where the displayed sections generate \(mB\) everywhere. Constants
and finite maxima preserve \(F_B\), after clearing degrees. Its zero
function is obtained by giving a generating family zero weights: at a
valuation-ring center every coefficient has norm at most one, and one
generator is a unit there.

Let \(\mathcal D_B\) consist of generically finite decreasing-net limits
from \(F_B\). Suppose \(\psi_i\downarrow\psi\) and
\(\psi_i=\inf_j f_{i,j}\) with \(f_{i,j}\in F_B\). Order triples
\((i,j,\epsilon)\) by increasing \(i\) and decreasing
\(f_{i,j}+\epsilon\), with \(\epsilon>0\). For two old triples, choose
\(i_3\) above both and \(\epsilon_3\) smaller than both old positive
errors. Then \[
\psi_{i_3}+\epsilon_3<
\min(f_{i_1,j_1}+\epsilon_1,f_{i_2,j_2}+\epsilon_2)
\] at every point. The right side is finite, so the inequality is still
strict where \(\psi_{i_3}=-\infty\). The exponentials of the inner
approximants minus the exponential of the finite right side decrease to
a strictly negative function. Their negative open sets cover the compact
analytic space; a finite subcover and directedness give one inner index
that works everywhere. Hence the triples are directed. The resulting net
\(f_{i,j}+\epsilon\) decreases to \(\psi\); at a pole use arbitrary
finite thresholds instead of subtracting infinities.

This proves the needed fixed-bundle nested-limit stability with
\textbf{finite inner approximants}. It gives no flattening theorem for
the common-pole original families after \(D\) is added. A single section
logarithm on \(B\) is in \(\mathcal D_B\): its floors are obtained by
adjoining a sufficiently divisible globally generating family of weight
\(-j\).

The fixed-bundle ample approximation granted by G0 therefore permits a
single finite \(F_A\) net for each genuine
\(\psi\in\operatorname{PSH}(Y,A)\). No countable extraction is required.

\subsection{2.3. Divided nets need not
decrease}\label{divided-nets-need-not-decrease}

Already on \(\mathbb P^1\), let \(\sigma\leq0\) be the logarithm of a
point-divisor section and consider \[
h_i=(1-i^{-2})\sigma,\qquad t_i=1-i^{-1}\quad(i\geq2).
\] The \(h_i\) decrease and \(h_i\leq t_i\sigma\), but their literal
quotients are \[h_i-t_i\sigma=(i^{-1}-i^{-2})\sigma.\] The positive
coefficient decreases for \(i\geq2\), so the quotients increase wherever
\(\sigma<0\). Their bundle classes tend to the ample remainder. Thus a
quotient-net monotonicity assertion cannot replace ample compactness.
Section 5 uses compactness instead.

\section{3. Existence of the retained algebraic
datum}\label{existence-of-the-retained-algebraic-datum}

\subsection{3.1. A smooth projective modification over every field in
the
domain}\label{a-smooth-projective-modification-over-every-field-in-the-domain}

Chow's lemma, in the form Stacks 0200 and 0201, gives an integral proper
birational refinement \(Z\to X\) with an immersion
\(Z\to\mathbb P^N_k\). Take the reduced closure of the common dense
open. Since \(Z\) is proper over \(k\), the immersion is closed, so
\(Z\) is projective and integral.

An integral finite-type scheme over any characteristic-zero field is
Noetherian, excellent, generically reduced, and a scheme over
\(\operatorname{Spec}\mathbb Q\). It need not be of finite type over
\(\mathbb Q\). The exact scope of Temkin's Theorem 1.2.1,
arXiv:0904.1592v2, applies to this \(Z\) and supplies a finite
regular-final blowup sequence. Integral blowup charts are subrings of
the common function field. Each blowup is projective by its degree-one
generated Rees algebra; composing the finite sequence retains
projectivity and birationality. The final \(Y\) is integral, regular,
projective, and of finite type over \(k\). Characteristic zero makes
\(k\) perfect, and Stacks 0B8X identifies regularity with smoothness in
this domain. Set \(\pi:Y\to X\) to be the composite.

Stacks 0AY8 applies with normal integral base \(X\), reduced integral
source \(Y\), and generic fiber \(\operatorname{Spec}k(X)\). It gives
\[\pi_*\mathcal O_Y=\mathcal O_X.                         \tag{3.1}\]
Locally trivialize any original bundle and apply (3.1) on that original
open. Gluing then gives, for each integral original bundle \(F\),
\[H^0(Y,\pi^*F)=H^0(X,F),                              \tag{3.2}\] with
the actual pullback map, not an abstract equality of dimensions.

\subsection{3.2. Big actual bundles exist on every original
carrier}\label{big-actual-bundles-exist-on-every-original-carrier}

Choose an effective very ample Cartier divisor \(H\) on \(Y\) containing
none of the finitely many exceptional prime divisors of \(\pi\). The
infinite field \(k\) permits a section avoiding this finite set of
linear conditions. Its pushforward \(B=\pi_*H\) is an effective Weil
divisor on \(X\). Smooth local rings are regular local UFDs (Stacks
0AG0), so \(B\) is an actual Cartier divisor. Birational equality away
from the exceptional locus gives \[\pi^*B=H+E,\] where \(E\) is
effective and exceptional: the pullback of an effective Cartier divisor
is effective, and \(H\) contains no exceptional prime. Multiplication by
\(s_E^m\) and (3.2) inject \[H^0(Y,mH)\hookrightarrow H^0(X,mB).\] The
left dimension grows with positive leading term of degree \(n\). Hence
\(\mathcal O_X(B)\) is an actual big line bundle. No ample bundle on the
original \(X\) has been introduced.

\subsection{3.3. Every specified big rational bundle has a polar
presentation}\label{every-specified-big-rational-bundle-has-a-polar-presentation}

Now fix any actual big rational \(L\) on \(X\) and clear denominators to
\(F=q\pi^*L\) on \(Y\). By (3.2), \(F\) has the same positive
section-growth limsup as \(qL\). Choose any actual very ample \(M\) on
\(Y\), and a nonzero hyperplane section \(T\in|M|\). We do not require
\(T\) to be reduced or smooth.

On the projective scheme \(T\), choose an ample line bundle \(N\)
sufficiently positive that \(N-F|_T\) is globally generated. A section
of that difference avoiding all the finitely many associated points is
regular. Such a section exists over the infinite field \(k\), since the
prohibited sections form a finite union of proper linear subspaces. Its
powers inject \(F^m|_T\) into \(N^m\). Serre vanishing and the coherent
Euler-polynomial support bound give
\[h^0(T,mF|_T)\leq h^0(T,mN)=O(m^{n-1}).\] The restriction sequence for
the effective Cartier \(T\) implies
\[h^0(Y,mF-M)\geq h^0(Y,mF)-O(m^{n-1}).\] Along arbitrarily large
integers realizing the positive \(m^n\) limsup, this is positive. Fix
one such \(m\) and one actual nonzero section with divisor \(G\geq0\).
Then \[\pi^*L=(qm)^{-1}M+(qm)^{-1}G.\] Retain \(A=(qm)^{-1}M\),
\(D=(qm)^{-1}G\), and this actual isomorphism. This proves (1.1) for
every specified big rational \(L\), using no Bertini statement or finite
generation.

\subsection{3.4. Whole bad fibers, rather than selected exceptional
points}\label{whole-bad-fibers-rather-than-selected-exceptional-points}

Let \(U_\pi\) be the maximal algebraic open over which \(\pi\) is an
isomorphism, and put \(B_\pi=X\setminus U_\pi\). Stacks 0BFP implies
that \(B_\pi\) has codimension at least two. More particularly, if there
is a quasi-finite point over a normal-base point, that base point
belongs to \(U_\pi\). Thus every irreducible component of a geometric
bad fiber has positive dimension: a zero-dimensional component would
contain a quasi-finite point. This argument applies after field
extension, since the smooth base remains normal and the birational
morphism remains proper.

On a reduced integral component \(F_0\) of such a geometric fiber,
\(\pi^*L\) is literally trivial, while \(A|_{F_0}\) is ample.
Consequently \(\mathcal O(mD)|_{F_0}\simeq\mathcal O(-mA)|_{F_0}\) for a
cleared positive integer \(m\). An antiample line bundle on a
positive-dimensional proper integral variety has no nonzero section.
Indeed a nonzero section and its powers would inject the unbounded
spaces \(H^0(F_0,rmA)\) into the finite space \(H^0(F_0,\mathcal O)\), a
contradiction. The canonical divisor section therefore vanishes on every
such component. We have proved \[\pi^{-1}(B_\pi)\subset |D|.\]

Let \[P=\pi(|D|)\subset X.\] This is a proper closed set: properness
makes the image closed, and the original generic point has its generic
singleton preimage outside \(D\). It contains \(B_\pi\). Outside
\(B_\pi\) every fiber is a singleton, and over \(B_\pi\) the entire
fiber lies in \(D\). It follows that \[\pi^{-1}(P)=|D|,
\qquad Y\setminus|D|\simeq X\setminus P.                 \tag{3.3}\] The
identities hold as closed supports, which is all that seminorms and
measures use. Analytification of closed supports and of the open
isomorphism gives \[p^{-1}(P^{\rm an})=|D|^{\rm an},\quad
p:Y^{\rm an}\setminus|D|^{\rm an}\simeq X^{\rm an}\setminus P^{\rm an}, \tag{3.4}\]
where \(p=\pi^{\rm an}\).

\section{4. Analytic centers, the original reference, and actual section
descent}\label{analytic-centers-the-original-reference-and-actual-section-descent}

\subsection{\texorpdfstring{4.1. Compactness and quotient topology
without projectivity of
\(X\)}{4.1. Compactness and quotient topology without projectivity of X}}\label{compactness-and-quotient-topology-without-projectivity-of-x}

For an analytic point \(x\), properness extends its associated morphism
\(\operatorname{Spec}\mathcal H(x)\to X\) uniquely to its valuation
ring. This gives the center \(c_X(x)\). Centers are compatible with
morphisms of proper varieties: compose the extended valuation-ring
morphism and use uniqueness. They are compatible with isometric
residue-field extensions for the same reason. These are actual centers,
not chosen valuations above the point.

Separatedness of algebraic \(X\) gives Hausdorffness of \(X^{\rm an}\)
by Berkovich 1993, 2.6.1 and 2.6.7. Compactness also has a direct center
proof. Choose a finite affine cover \(U_i=\operatorname{Spec}R_i\) of
\(X\). The set of seminorms on \(R_i\) bounded by one on every element
is a closed subset of \([0,1]^{R_i}\), since the multiplicative and
ultrametric identities are closed conditions. It is compact by
Tychonoff, without a cardinality restriction. Its image in
\(X^{\rm an}\) is exactly \(c_X^{-1}(U_i)\): boundedness supplies the
valuation-ring map, and conversely a center in \(U_i\) makes every
regular coefficient bounded by one. Those compact sets form a finite
cover of \(X^{\rm an}\).

The analytic morphism \(p\) is continuous and surjective. For
surjectivity one can use Berkovich 2.6.8, whose proof is by extension of
field valuations. In this particular proper setting, base change to
\(\mathcal H(x)\) and choose a closed point of the nonempty algebraic
fiber; its finite residue extension has the extended complete-field norm
and produces an analytic point above \(x\). Since \(Y^{\rm an}\) is
compact and \(X^{\rm an}\) Hausdorff, \(p\) is closed and is a quotient
map. Its compact fibers are the actual analytic fibers. Its restriction
to the open complements in (3.4) is the analytification of an actual
algebraic isomorphism. This proof does not use the omitted proof of
Berkovich 2.6.9.

\subsection{4.2. Canonical trivial metrics on the original
carrier}\label{canonical-trivial-metrics-on-the-original-carrier}

For an actual line bundle \(F\) on proper \(X\), define its canonical
trivial metric by making an algebraic frame at the valuation-ring center
have norm one. On \(c_X^{-1}(U_i)\), a frame \(e_i\) of \(F|_{U_i}\)
therefore has \(\|e_i\|_0=1\). On overlaps transition functions are
units at the center, hence have seminorm one. The metrics glue.

This metric is continuous. For any other algebraic frame \(e\) on an
open \(U\), its norm on each compact center set meeting \(U^{\rm an}\)
is the continuous seminorm of the corresponding transition function
\(e/e_i\). These sets give a finite relatively closed cover of
\(U^{\rm an}\). The closed-cover gluing criterion proves continuity on
\(U^{\rm an}\). Center compatibility proves the literal tensor and
pullback identities for canonical metrics, at every analytic point.
Clearing powers supplies the same statements for rational bundles. This
reconstructs the extension of the canonical-reference construction used
in BJ 1.24 to the present proper original carrier.

Write \[\sigma=\log\|s_D\|_0\leq0\] on \(Y^{\rm an}\), clearing the
rational divisor. It is continuous into \(\mathbb R\cup\{-\infty\}\) and
is \(-\infty\) exactly on \(|D|^{\rm an}\). For any genuine
\(\psi\in\operatorname{PSH}(Y,A)\), the function \(\sigma+\psi\) is usc,
bounded above, and constant on every fiber of \(p\). It is \(-\infty\)
on the whole preimage of \(P\), and all remaining fibers are singletons.
Let \[F(\psi)=u,\qquad p^*u=\sigma+\psi.                    \tag{4.1}\]
Closedness and surjectivity give \(\{u\geq a\}=p\{\sigma+\psi\geq a\}\)
closed for every finite \(a\), so \(u\) is usc. Its all-point value on
\(P\) is fixed, not inferred from divisorial density. At the generic
trivial point its value is \(\psi(v_{\rm triv})\), hence it is
generically finite. Put \(\tau=F(0)\); its pole set is exactly
\(P^{\rm an}\).

\subsection{4.3. A single genuine original-section
net}\label{a-single-genuine-original-section-net}

Choose a finite fixed-\(A\) net \(\psi_\alpha\downarrow\psi\) from
Section 2.2. At each stage clear all section degrees and bundle
denominators to \(r>0\). Multiply each section of \(rA\) by the
canonical section \(s_{rD}\). They are actual sections of \(r\pi^*L\).
By (3.2), each is the pullback of an actual section of \(rL\) on the
original \(X\). Canonical metric compatibility makes their logarithmic
norm formulas equal at \textbf{every} analytic point, with the same
constants. Let \(u_\alpha\) be their original maximum. Then
\[p^*u_\alpha=\sigma+\psi_\alpha,\qquad u_\alpha\downarrow F(\psi) \tag{4.2}\]
pointwise on \(X^{\rm an}\), including all poles. Neither the original
maxima nor \(\tau\) need be finite at their common base locus. The
flattening was performed before adding \(\sigma\); no strict inequality
at an original common pole was used. The reference \(\tau\) itself is a
finite original section maximum, since zero is a finite FS function on
\(A\).

These are actual singular metrics on \(L\): in an original local frame
\(e\), their raw weights are \(u_\alpha-\log\|e\|_0\), and their
decreasing limit is \(u-\log\|e\|_0\). They are not abstract energy
completions or an arbitrarily assigned lift.

\section{5. A direct repair of the entire fixed-divisor
cone}\label{a-direct-repair-of-the-entire-fixed-divisor-cone}

This section reconstructs the nonample foundations proposed in
continuation114. They are used only on the projective smooth \(Y\).
Define a numerical psh function by the original single-net definition:
it is a generically finite decreasing-net limit of actual
real-coefficient, generically finite FS maxima \(\phi_i\) in rational
bundles \(L_i\), with \(c_1(L_i)\to\theta\). No nested varying-class
closure is included as a premise.

\subsection{5.1. Rounding and embedding, directly in the defining
net}\label{rounding-and-embedding-directly-in-the-defining-net}

For positive rational \(\delta\), round the finitely many coefficients
of \(\phi_i\) upward to rational coefficients with error at most its
degree times \(\delta\). The resulting \(f_{i,\delta}\) is in the
\textbf{same} bundle and satisfies
\[\phi_i\leq f_{i,\delta}\leq\phi_i+\delta.\] Set
\(g_{i,\delta}=f_{i,\delta}+\delta\). Order pairs by increasing \(i\)
and decreasing \(g_{i,\delta}\). Given two pairs, take \(i_3\) above
both and \(2\delta_3\leq\min(\delta_1,\delta_2)\). Then
\[g_{i_3,\delta_3}\leq\phi_{i_3}+2\delta_3\leq
\min(\phi_{i_1}+\delta_1,\phi_{i_2}+\delta_2)\leq
\min(g_{i_1,\delta_1},g_{i_2,\delta_2}).\] These non-strict inequalities
hold at common poles as well. The pair set is directed; its projection
to the original index is final. Approximation at each finite value, or
below every real threshold at a pole, proves
\(g_{i,\delta}\downarrow\phi\) and preserves class convergence.

Choose an actual ample rational \(H\) with \(c_1(H)-\theta\) ample.
Eventually \(H-L_i\) is an actual ample rational bundle. Multiplication
by a zero-weight generating family of a divisible power of this
difference realizes \(\phi_i\) as an actual \(H\)-section maximum. If it
has base poles, adjoin a zero reference floor of weight \(-m\). The
product index \((i,m)\) gives a single decreasing net of finite
fixed-\(H\) maxima to \(\phi\). Thus \(\phi\) belongs to the genuine
actual-ample \(H\) cone. The ample G0 divisorial and usc-obstacle
principles now apply to it. A finite list of numerical functions can be
put in one common \(H\). This argument does not invoke nonample
regularization or varying-class closure.

\subsection{5.2. Homogeneous projection by a finite rational
polytope}\label{homogeneous-projection-by-a-finite-rational-polytope}

For \(\phi(v_{\rm triv})\geq0\) put
\[\widehat\phi(v)=\inf_{t>0}\frac{\phi(tv)}t.\] The two infima commute:
\[\inf_i\widehat\phi_i(v)=\inf_t\inf_i\phi_i(tv)/t=\widehat\phi(v).\]
For a rational finite FS expression \(m^{-1}\max_j(x_j+\lambda_j)\), set
\[\mathscr P=\{w_j\geq0:\ \sum w_j=1,\ \sum w_j\lambda_j\geq0\}.\] Then
\[\inf_{r\geq0}\max_j(x_j+r\lambda_j)
=\max_{w\in\operatorname{Vert}(\mathscr P)}\sum_jw_jx_j.       \tag{5.1}\]
Here is a direct one-dimensional proof. Every feasible weighted
expression bounds the left side from below. For finite \(x_j\), the
convex piecewise-linear function on the left attains its minimum at
zero, at a breakpoint with active slopes of opposite signs, or as a
limit at infinity. At a zero minimum an active nonnegative slope
supplies a unit-vector vertex. At a breakpoint, a zero weighted average
of two active opposite slopes supplies a vertex. In the final case the
maximal slope is zero and an active zero-slope unit vector supplies the
limiting value. These witnesses give equality. Vertices are the
indicated unit vectors and two-slope zero averages, hence rational. For
\(x_j=-\infty\), replace them by \(\max(x_j,-N)\) and commute the infima
over \(N\) and \(r\); the fixed finite maximum of vertex expressions
passes to its decreasing limit. Zero coefficients omit those coordinates
from a monomial.

Clear a common denominator \(b\) of the vertex weights. The products
\(\prod_j s_j^{bw_j}\) are actual sections of \(mbL\). Their homogeneous
logarithmic maximum divided by \(mb\) is \(m^{-1}\) times the right side
of (5.1), with all-point values including base poles. Hence
\(\widehat\phi_i\) are actual homogeneous rational FS functions in the
same bundles as \(\phi_i\), with generic value zero. They decrease to
\(\widehat\phi\), so the latter is genuine by the single defining net.
If \(h\) is already homogeneous, \(\widehat h=h\).

\subsection{5.3. Pointwise homogeneous decomposition by
convexity}\label{pointwise-homogeneous-decomposition-by-convexity}

Let \(C=\phi(v_{\rm triv})\) and \(h_\lambda=\widehat{\phi-\lambda}\)
for \(\lambda<C\). A tail of the defining net has nonnegative generic
values after subtracting \(\lambda\), so Section 5.2 applies. It gives
genuine homogeneous \(h_\lambda\), and \(h_\lambda+\lambda\leq\phi\) by
taking \(t=1\).

For a fixed actual point \(v\), the decreasing limit \(g(t)=\phi(tv)\)
is convex, decreasing, and bounded above by \(C\). It is either finite
on \((0,\infty)\) or identically \(-\infty\) there. To see the latter
assertion, a \(-\infty\) value at \(t_0\) and the common finite upper
bound at a smaller \(s\) force the limit to be \(-\infty\) at all
\(t\in(s,t_0)\) by convexity; monotonicity handles larger \(t\). Let
\(s\downarrow0\).

In the finite case choose a supporting line \(at+\lambda_0\) at \(t=1\).
Its slope is nonpositive, and \(\lambda_0\leq C\) by letting
\(t\downarrow0\) in \(at+\lambda_0\leq g(t)\leq C\). If \(\lambda_0<C\),
then \(h_{\lambda_0}(v)=a\), giving equality at \(v\). If
\(\lambda_0=C\), lower the intercept by \(\epsilon\):
\(h_{C-\epsilon}(v)\geq a\), giving the desired value within
\(\epsilon\). The identically infinite case is immediate. Therefore at
every actual point
\[\phi=\sup_{\lambda<C}(h_\lambda+\lambda).                \tag{5.2}\]
No usc replacement was made in this equality.

\subsection{5.4. Literal divisor division and ample
compactness}\label{literal-divisor-division-and-ample-compactness}

Let \(B=A+D\) be the retained actual bundle on \(Y\). Suppose \(h\) is
homogeneous numerical \(B\)-psh and \(h\leq\sigma\). Its homogeneous
actual FS approximants \(h_i\) from Section 5.2 decrease in bundles
\(B_i\) with classes tending to \(c_1(B)\). If \(D\ne0\), the level set
\(K=\{\sigma=-1\}\) is compact. For \(0<t<1\) rational, the open sets
\(\{h_i<-t\}\) increase to a cover of \(K\), so a tail satisfies that
inequality throughout \(K\). Homogeneity then gives \(h_i\leq t\sigma\)
at every function-field divisorial valuation with \(\sigma<0\), by
scaling to \(K\); when \(\sigma=0\), use \(h_i\leq0\). The common-ample
embedding and its usc-obstacle principle extend the inequality to all
points.

Use admissible pairs \((i,t)\) with \(i\) increasing and \(t\uparrow1\).
They are directed by first choosing a larger \(t\) and then a
sufficiently late \(i\). For each displayed section, the divisorial
inequality says its order along every prime divisor of \(D\) is at least
its degree times \(tD\). Clearing denominators and normality of \(Y\)
make division by \(s_{ptD}\) an actual global section of \(p(B_i-tD)\).
The resulting actual FS function \(r_{i,t}\) is the literal quotient on
the divisor complement and its actual extended value elsewhere; it has
generic value zero.

Refine by rational \(c>1\), tending to one, with \(cA-(B_i-tD)\) ample.
Openness of the ample cone gives admissible triples arbitrarily late and
with \(t,c\) arbitrarily close to one. Direct their indices by
increasing \(i,t\) and decreasing \(c\); a common upper triple is
obtained by fixing a smaller admissible \(c\) and taking the other
indices sufficiently late. Multiplying zero-weight generators of
\(A-(B_i-tD)/c\) realizes \(r_{i,t}/c\) in the actual \(A\) cone. If its
maxima have poles, the fixed-bundle floors from Section 5.1 give genuine
ample membership.

The normalized actual \(A\) cone is compact by G0. A convergent subnet
has a genuine quotient \(r\). At each divisorial point its values
converge explicitly to \(h-\sigma\), since \(h_i\), \(t\), and \(c\)
have the stated limits. Multiply a finite fixed-\(A\) net to \(r\) by
\(s_D\); this proves \(\sigma+r\) is genuine numerical \(B\)-psh. Both
it and \(h\) embed in one actual ample cone. Ample divisorial
determination now gives \(h=\sigma+r\) at every point. The quotient net
itself was not claimed decreasing. When \(D=0\), the same proof uses
\(h_i\) directly and omits \(K\) and division.

For an arbitrary numerical \(B\)-psh \(v\leq\sigma+C_0\), (5.2) gives
\(h_\lambda+\lambda\leq v\). Scaling any divisorial point and letting
the scale tend to infinity gives \(h_\lambda\leq\sigma\) there, hence
everywhere by the ample obstacle principle. Divide to get normalized
\(r_\lambda\). The family \(r_\lambda+\lambda\) is bounded above by
\(C=v(v_{\rm triv})\). Its ample usc envelope \(r\) is genuine. The
ample divisorial-nonnegligibility principle preserves its divisorial
values; thus (5.2) gives \(v=\sigma+r\) on divisorial points. Put both
sides in one ample cone and apply determination again. We have proved
the full all-point cone identity
\[\{v\in\operatorname{PSH}_{\rm numerical}(Y,B):v\leq\sigma+O(1)\}
=\sigma+\operatorname{PSH}(Y,A).                         \tag{5.3}\] The
quotient is unique, since \(\sigma\) is finite on function-field
divisorial points. Only actual-ample compactness, envelopes, and
determination were used; the nonample ingredients were reconstructed
from one defining FS net.

\subsection{5.5. Original section cone and its genuine
inverse}\label{original-section-cone-and-its-genuine-inverse}

Let \(\operatorname{PSH}_{\rm sec}(X,L)\) be the generically finite
single decreasing-net closure of actual original \(L\)-section maxima,
allowing their common base poles. Pulling its defining net to \(Y\)
gives a single defining numerical \(B\) net. Define the retained cone on
the original carrier by \[\mathcal P_{\mathcal M}
=\{u\in\operatorname{PSH}_{\rm sec}(X,L):u\leq\tau+C\text{ for some }C\in\mathbb R\}.\]
Section 4 proves
\(F(\operatorname{PSH}(Y,A))\subset\mathcal P_{\mathcal M}\). Conversely
the actual pullback of \(u\in\mathcal P_{\mathcal M}\) satisfies
\(p^*u\leq\sigma+C\), so (5.3) supplies its unique genuine quotient
\(Q(u)\). Surjectivity of \(p\) gives \(u=F(Q(u))\) at every point. Thus
\[F:\operatorname{PSH}(Y,A)\simeq\mathcal P_{\mathcal M},\qquad Q=F^{-1}. \tag{5.4}\]
This is genuine section descent, not an arbitrarily chosen lift. It
reflects order: compare on divisorial points, cancel their finite
\(\sigma\), and use the ample usc-obstacle principle. It commutes with
constants and finite maxima. Its divisorial topology is the original
pointwise topology on the common function field, translated by the fixed
finite values of \(\sigma\). Hence (5.4) is a homeomorphism, preserving
suprema at the generic trivial point. The normalized original cone is
compact; decreasing original limits have a genuine quotient and an
actual original section net by Section 4.

The numerical zero-class repair in continuation114 is also valid within
these foundations. Given a single defining FS net with classes tending
to zero, generic values \(C_i\downarrow C\), refine by
\(\epsilon\downarrow0\) with \(\epsilon A-L_i\) ample. The functions
\((\phi_i-C_i)/\epsilon\) belong to the normalized actual \(A\) cone. At
every divisorial point \(v\), compactness and continuous evaluation give
a common finite lower bound \(M_v\). Thus
\[\epsilon M_v\leq\phi_i(v)-C_i\leq0.\] The limit equals \(C\)
divisorially and, after one ample embedding, everywhere. This
establishes a constants-only zero class without applying varying-class
closure to \(m\phi\). It is not used to extend the present task to
zero-volume nonample classes.

\section{6. Original finite energy, fiber maxima, and a finite-energy
pole}\label{original-finite-energy-fiber-maxima-and-a-finite-energy-pole}

Define, intrinsically for the retained genuine cone,
\[E_{\mathcal M}(u)=E_A(Q(u)),\qquad
\mathcal E^1_{\mathcal M}=\{u\in\mathcal P_{\mathcal M}:E_{\mathcal M}(u)>-\infty\}. \tag{6.1}\]
This energy is attached to actual original metrics, with its genuine
quotient uniquely recovered in (5.4). It is concave, monotone, weakly
usc, and decreasing-net continuous by G0 and the order homeomorphism.
Its constant law is \(E_{\mathcal M}(u+c)=E_{\mathcal M}(u)+Vc\).

For \(u=F(\psi)\) set
\[q_u(x)=\max_{y\in p^{-1}(x)}\psi(y).                    \tag{6.2}\]
The maximum is attained on every compact fiber, including value
\(-\infty\). The sets \(\{q_u\geq a\}=p\{\psi\geq a\}\) are closed.
Hence \(q_u\) is an actual upper bounded usc function on the original
analytic carrier. On its open complement of \(P\) it equals \(u-\tau\).
On \(P\), formula (6.2) gives a uniquely specified value from the
retained genuine quotient; it is never an expression
\(-\infty-(-\infty)\). If \(\psi_i\downarrow\psi\), then
\(q_{F(\psi_i)}\downarrow q_{F(\psi)}\): on a fixed fiber the closed
sets \(\{\psi_i\geq a\}\) decrease, and compactness excludes a
persistent level exceeding the limiting maximum. Constants and order are
respected. Additivity on bad fibers is not asserted or needed.

There is an actual \(w\in\mathcal E^1(Y,A)\) equal to \(-\infty\) on all
of \(D\). Choose \(r\) sufficiently large that \(rA-D\) has an effective
nonzero section. Multiplying by the divisor section gives a section of
\(rA\) vanishing along \(D\). Its normalized logarithm \(a\leq0\) is
genuine \(A\)-psh by fixed-bundle floors. Set
\[\alpha=\frac1{n+2},\qquad w=1-(1-a)^\alpha.             \tag{6.3}\]
The convex increasing function in (6.3) has slopes in \([0,1]\), value
zero at zero, and value \(-\infty\) at \(-\infty\). Its affine supports
applied to \(a\) are genuine \(A\) potentials by convexity with zero.
Their bounded-above ample envelope is exactly the usc composition \(w\),
proving genuine membership without nonample composition closure.

Here is the finite-energy estimate, including its necessary exponent. In
probability calibration let
\[\operatorname{Cap}(E)=\sup_{-1\leq b\leq0,\ b\in\operatorname{PSH}(A)}\operatorname{MA}_{A,\rm prob}(b)(E).\]
For every normalized genuine \(a\),
\(\operatorname{Cap}\{a\leq-s\}\leq n/s\). To justify this without
unbounded integration assumptions, apply signed integration by parts to
the bounded floors of \(a\). Telescope the difference between the
measure of \(b\) and that of zero. Each of the \(n\) terms is an
integral of \(b\in[-1,0]\) against the difference of two positive mixed
products of the same probability mass, and has absolute value at most
one. The zero-potential measure is the generic trivial Dirac measure,
where every normalized floor vanishes. Monotone convergence in the fixed
measure of \(b\) gives
\(\int(-a)\operatorname{MA}_{\rm prob}(b)\leq n\). Chebyshev and the
supremum over \(b\) prove the capacity bound.

For any normalized genuine \(z\), write \(z_s=\max(z,-s)\),
\(\mu_s=\operatorname{MA}_{\rm prob}(z_s)\). Energy concavity gives
\(J(z_s)\leq\int(-z_s)d\mu_s\). Locality and equal masses show, for
\(0<t\leq s\), that \(\mu_s\{z\leq-t\}=\mu_t\{z\leq-t\}\). For
\(t\geq1\), positivity of the mixed expansion gives
\(\operatorname{MA}_{\rm prob}(z_t/t)\geq t^{-n}\mu_t\). Layer cake
therefore gives
\[J(z_s)\leq1+\int_1^\infty t^n\operatorname{Cap}\{z\leq-t\}\,dt. \tag{6.4}\]
The factor \(t^n\) is retained; it is missing in a printed intermediate
line in BJ's proof. In (6.3),
\(\{w\leq-t\}=\{a\leq1-(1+t)^{1/\alpha}\}\), so its capacity is
\(O(t^{-1/\alpha})\). The integrand in (6.4) is \(O(t^{-2})\). The
energies of the floors are bounded below uniformly, and decreasing
energy continuity proves \(w\in\mathcal E^1(Y,A)\). We make no claim for
the larger printed exponent interval in BJ 7.48.

Since \(p^{-1}(P)=D\), (6.2) gives \(q_{F(w)}=-\infty\) on all of \(P\).
Every mixed finite-energy product on \(Y\) integrates \(w\) by G0, hence
gives \(D\) zero mass. This establishes the exceptional-set conclusion
used below before any target is chosen.

\section{7. Radon inversion and the original intrinsic target
class}\label{radon-inversion-and-the-original-intrinsic-target-class}

We first prove the relevant measure theorem. Suppose \(p:S\to T\) is a
closed continuous surjection of compact Hausdorff spaces, \(C\subset T\)
is closed, \(p^{-1}(C)=F\), and \(p:S\setminus F\simeq T\setminus C\)
has inverse \(j\). Every Radon \(\mu\) on \(T\) with \(\mu(C)=0\) has a
unique Radon inverse \(\nu\) on \(S\) satisfying \(\nu(F)=0\) and
\(p_*\nu=\mu\).

For \(g\in C(S)\) define
\[\Lambda(g)=\int_{T\setminus C}g(j(x))\,d\mu(x).\] This is a positive
bounded functional, and Riesz gives a Radon \(\nu\). To prove zero mass
on \(F\), choose compact \(K\subset T\setminus C\) with
\(\mu(T\setminus K)\) arbitrarily small. A continuous \(\rho\in[0,1]\)
equal to one on \(C\) and zero on \(K\) satisfies
\(\nu(F)\leq\Lambda(\rho\circ p)=\int\rho d\mu\), which is arbitrarily
small. Continuous tests pulled from \(T\) prove pushforward. Uniqueness
follows from the isomorphism off \(F\), zero mass on \(F\), and
regularity; there is no measurable-selection or arbitrary-fiber choice.
Integrals of upper bounded usc functions agree, including \(-\infty\),
by their continuous decreasing majorants or ordinary Borel transport.

Apply this to \(S=Y^{\rm an}\), \(T=X^{\rm an}\), \(C=P^{\rm an}\), and
\(F=D^{\rm an}\). For a positive Radon \(\mu\) of mass \(V\), define the
original retained dual energy
\[D_{\mathcal M}(\mu)=\sup_{u\in\mathcal E^1_{\mathcal M}}
\left(E_{\mathcal M}(u)-\int q_u\,d\mu\right).             \tag{7.1}\]
The quotient integral is relative to the actual reference; absolute
integrability of \(u\) or \(\tau\) is not asserted. If (7.1) is finite,
every \(q_u\) is integrable: otherwise its individual expression would
be \(+\infty\). In particular the explicit pole \(q_{F(w)}\) forces
\(\mu(P)=0\). The unique Radon inverse \(\nu\) exists and
\[\int q_u\,d\mu=\int Q(u)\,d\nu,\qquad
D_{\mathcal M}(\mu)=E_A^\vee(\nu).                       \tag{7.2}\] G0
proves that (7.1) is finite if and only if \textbf{every} \(q_u\),
\(u\in\mathcal E^1_{\mathcal M}\), is \(\mu\)-integrable. In the reverse
implication, the pole again excludes \(P\) first, then G0 is applied
through (7.2). Thus this target condition is intrinsic to the retained
original metric cone and its reference. It specifies the exact mass and
energy and imposes no toric, skeletal, or countability restriction. Zero
mass on \(P\) is a proved consequence, rather than a silent support
assumption.

For completeness, zero-exceptional Radon inversion is weakly continuous
at zero-exceptional limits. If \(\mu_i\to\mu\) weakly, all masses equal
\(V\), and all \(\mu_i(P)=\mu(P)=0\), let \(\nu_i,\nu\) be their
inverses. For \(g\in C(Y)\) choose a bump \(\rho=1\) near \(P\) with
\(\int\rho\,d\mu\) small. The function \((1-\rho\circ p)g\) descends off
\(P\) and extends continuously by zero across \(P\). Its integrals
converge. The remaining integrals have absolute values at most
\(\|g\|_\infty\int\rho\,d\mu_i\), which are eventually small. This
proves weak convergence of the inverses for arbitrary nets.

\section{8. Pretarget mixed pairing and the local
progress}\label{pretarget-mixed-pairing-and-the-local-progress}

\subsection{8.1. A measure-valued pairing, without calling it the
missing
operator}\label{a-measure-valued-pairing-without-calling-it-the-missing-operator}

For \(u_i=F(\psi_i)\in\mathcal E^1_{\mathcal M}\) define
\[\widetilde M_{\mathcal M}(u_1,\ldots,u_n)
=p_*\bigwedge_{i=1}^n(c_1(A)+dd^c\psi_i).                \tag{8.1}\]
This is constructed for the entire cone before any target is chosen. It
is a positive symmetric Radon measure of raw mass \(V\). It gives \(P\)
zero mass by Section 6. Its local finite and limiting properties proved
next give more information than a bare transported-measure definition,
but we retain the tilde until the remaining bridge in Section 10 is
established.

G0 and compact pushforward give unbounded decreasing-net weak
convergence in every entry, with genuine finite-energy limits.
Integrable weighted limits also transport. The defining identity has the
sign \[\int\psi_0\,\bigwedge_{i=1}^n(c_1(A)+dd^c\psi_i)
=(c_1(A),\psi_0)\prod_{i=1}^n(c_1(A),\psi_i)
-(c_1(A),0)\prod_{i=1}^n(c_1(A),\psi_i).                 \tag{8.2}\] It
is this difference, not its reverse, that is passed to the limit. On the
original carrier, (8.2) gives
\(q_{u_0}\widetilde M_{\mathcal M}=p_*(\psi_0 M_A)\). For integrable
decreasing tests and decreasing finite-energy mixed factors, G0 gives
the weighted weak limit. A difference of two such tests is treated term
by term; the difference is not falsely called decreasing. Ordinary weak
pushforward is legitimate on the compact spaces established in Section
4.

\subsection{8.2. Literal local-frame signs at continuous section
stages}\label{literal-local-frame-signs-at-continuous-section-stages}

Let \(W\subset X^{\rm an}\setminus P^{\rm an}\) be an ordinary open
frame neighborhood with compact closure still inside a common frame
region. Pull it to \(Y\setminus D\) using the actual isomorphism (3.4).
Choose compatible frames \(e_A,e_D\) under \(\iota\), so the pulled
original frame is \(e_L=e_Ae_D\). Write \(s_D=f e_D\), where \(f\) is an
algebraic, hence analytic, unit on this region. Let
\(h_A=-\log\|e_A\|_0\), \(h_D=-\log\|e_D\|_0\). Then
\[\sigma=\log|f|-h_D,\qquad
u+h_L=\psi+h_A+\log|f|.                                \tag{8.3}\] The
unit logarithm has a \textbf{plus} sign. On a compact frame neighborhood
its value and the reference weights are finite and bounded.

For a finite \(F_A\) section maximum \(\psi_\alpha\), (4.2) gives the
actual original section maximum, with raw weight
\(\psi_\alpha+h_A+\log|f|\). Its finite-stage local metric intersection
product is the usual product for that actual section metric. BE 7.14
constructs smooth semipositive regularizations of the finite \(A\)
metric by regularized maxima. BE 8.4 gives the uniform-to-weak
continuous product and its identification with the genuine finite
intersection product. CLD 18.1.2 gives the zero Hessian of \(\log|f|\),
and CLD 8.8.5 gives its invariance at continuous local stages.
Restriction through the analytic isomorphism therefore identifies the
finite \textbf{actual original local} product with the finite ample
product on \(p^{-1}(W)\). This is a pretarget metric identity, not an
identity imposed by a desired measure.

We use those primary statements only in their continuous domains. They
do not establish a bounded discontinuous local product merely by taking
uniform limits.

\subsection{8.3. A precise decreasing local section
category}\label{a-precise-decreasing-local-section-category}

Here is the part that can be proved directly from G0. On the preceding
\(W\), consider nets of raw weights
\[\lambda_{i,\alpha}=\psi_{i,\alpha}+h_A+\log|f|,\] where
\(\psi_{i,\alpha}\) are finite actual \(F_A\) maxima, decreasing on all
of \(Y^{\rm an}\) to a genuine bounded \(A\)-psh endpoint \(\psi_i\).
The common directed index is obtained by a finite product of index sets.
Call these \emph{divided global-section nets on \(W\)}. They are actual
original-section local weights by (4.2).

For \(g\in C_c(W)\), pull \(g\) through the open isomorphism and extend
it by zero on \(Y\). G0's bounded decreasing-net continuity proves
convergence of its integrals against the finite local intersection
products in Section 8.2. The resulting local Radon measure is
\[M_A(\psi_1,\ldots,\psi_n)|_{p^{-1}(W)},                \tag{8.4}\]
read on \(W\) through the isomorphism. This conclusion applies even when
every \(\psi_i\) is bounded but discontinuous. It uses monotone ample
product convergence, not uniform convergence to the endpoint.

The limit is independent of the particular divided global-section nets.
If two bounded global endpoint tuples have equal raw weights on an
original open \(W\), their fixed reference and unit terms in (8.3)
cancel, leaving equal ample scalar potentials there. G0's mixed locality
makes (8.4) equal on \(W\). This also proves independence from changes
of the global endpoints outside \(W\), within this category. Frame
changes are handled at \textbf{each} finite stage: changing the original
frame subtracts the logarithm of its analytic unit from the raw weight,
without changing the finite local product. The same transformed
approximants pass through the decreasing construction, so their limits
agree as well. The fixed unit term in (8.3) has thus passed through the
bounded construction; its zero-Hessian property was not simply asserted
at the discontinuous endpoint.

One can also see the requisite ample locality directly from the granted
strict-winning identity. For equal potentials \(a,b\) on \(W\), apply
strict-winning locality to \(a+\epsilon,b\) there, then let
\(\epsilon\downarrow0\) in the genuine decreasing products. Exchange
\(a,b\) and polarize. All entries remain in the actual ample G0 domain.
Weighted versions use (8.2) with two separate decreasing terms.

This proves a genuine decreasing local limit for a specified
actual-section category, with all-point endpoints and germ independence
inside that category. It does not yet prove the comparison with a
construction made from only locally defined contact nets and a collar;
that comparison is the remaining issue.

\subsection{8.4. Relative floors and an unbounded candidate in the same
category}\label{relative-floors-and-an-unbounded-candidate-in-the-same-category}

For \(u_i=F(\psi_i)\in\mathcal E^1_{\mathcal M}\) put
\[u_{i,j}=\max(u_i,\tau-j)=F(\max(\psi_i,-j)).             \tag{8.5}\]
These are genuine bounded-quotient metrics. On compact original frame
neighborhoods off \(P\), their raw weights are bounded, possibly
discontinuous, and belong to Section 8.3. Absolute original floors
\(\max(u_i,-j)\) are not substituted for (8.5).

For the mixed product let \(G_j=\bigcap_i\{\psi_i>-j\}\). G0's mixed
strict-winning locality, obtained by polarization, gives
\[\mathbf1_{G_j}M_A(\max(\psi_1,-j),\ldots,\max(\psi_n,-j))
=\mathbf1_{G_j}M_A(\psi_1,\ldots,\psi_n).                \tag{8.6}\]
Each \(\psi_i\) is integrable for the latter measure, so its \(-\infty\)
set has zero mass. The measure also gives \(D\) zero mass. The
restrictions of the left side to \(Y\setminus D\) therefore increase to
the full genuine ample mixed measure. Their original restrictions give a
Radon candidate of mass exactly \(V\), equal to (8.1), and zero
extension across \(P\) is unique for a measure required to charge that
pole set by zero.

The restricted limit is independent of floor levels tending to infinity.
More generally, a bounded genuine \(A\)-psh reference \(b\) can replace
zero in \(\max(\psi_i,b-j)\); the same strict-winning and integrability
argument gives the same restriction. Weighted limits for integrable
genuine quotient tests follow from (8.2) and G0. Fixed frame-unit
transformations continue to agree at the bounded-floor stages and hence
in (8.6).

These assertions prove the relative nonpluripolar \emph{candidate} in
the divided category, including mass and weighted limits. They do not
silently define an independent local Bedford--Taylor product by
declaring (8.1) to be it.

\section{9. Original envelopes, all continuous variations, and
conditional
equations}\label{original-envelopes-all-continuous-variations-and-conditional-equations}

The deductions of this section apply to the pretarget pairing
\(\widetilde M\). They become an actual P16.2 equation theorem if the
local bridge in Section 10 is established.

\subsection{9.1. Genuine minorants and all original continuous majorants
agree}\label{genuine-minorants-and-all-original-continuous-majorants-agree}

For an upper bounded original obstacle \(h\), let
\[\mathscr O_{\mathcal M}(h)=\sup\{E_{\mathcal M}(v):v\in\mathcal P_{\mathcal M},\ v\leq h\}.\]
Take \(u=F(\psi)\in\mathcal E^1_{\mathcal M}\) and
\(f\in C^0(X^{\rm an})\). Cancel finite \(\sigma\) at function-field
divisorial points and apply the actual ample usc-obstacle principle.
This proves the exact equivalence
\[F(\eta)\leq u+tf\text{ on }X^{\rm an}\quad\Longleftrightarrow\quad
\eta\leq w_t:=\psi+t p^*f\text{ on }Y^{\rm an}.           \tag{9.1}\]
The pole fibers give no lost condition, because the divisorial
inequalities determine the genuine ample potentials everywhere. The
ample envelope \(P_A(w_t)\) lies between \(\psi\pm |t|\|f\|_\infty\),
hence is actual finite energy. Thus
\[\mathscr O_{\mathcal M}(u+tf)=E_A(P_A(w_t))=E_A^\downarrow(w_t). \tag{9.2}\]

There is also an outer definition using \textbf{all original continuous
relative majorants}: \[\mathscr O^{\rm out}_{\mathcal M}(u+tf)
=\inf_{b\in C^0(X^{\rm an}),\ q_u+tf\leq b}
E_{\mathcal M}(F(P_A(p^*b))).                            \tag{9.3}\] For
any \(b\in C^0(X)\), (9.1) and the fiber maximum show
\[u+tf\leq\tau+b\quad\Longleftrightarrow\quad q_u+tf\leq b.\] Compact
Hausdorff normality supplies continuous majorants of any upper bounded
usc function whose pointwise infimum is that function; finite minima
make the family a decreasing net. Let \(b_i\downarrow q_u+tf\) be this
net of all original majorants, and put \(\theta_i=P_A(p^*b_i)\). These
are genuine and decrease, bounded below by \(P_A(w_t)\) and above by a
fixed continuous bound. G0 gives a genuine finite-energy limit
\(\theta\).

Every function-field divisorial point has a unique analytic birational
lift: apply the proper valuative criterion to its common function-field
valuation ring. Its support is the original generic point, so \(p\) has
a singleton analytic fiber there and \(q_u(p(v))=\psi(v)\). It follows
that \(\theta(v)\leq w_t(v)\) for every such divisorial point. Ample usc
determination extends this inequality to all \(Y\). Conversely
\(\theta_i\geq P_A(w_t)\), since \(p^*b_i\geq w_t\). Hence
\(\theta=P_A(w_t)\). Decreasing energy continuity gives
\[\mathscr O^{\rm out}_{\mathcal M}(u+tf)=\mathscr O_{\mathcal M}(u+tf). \tag{9.4}\]
Only original continuous majorants were used in (9.3); no contraction to
a projective quotient restricts the variations. Continuous relative
obstacles \(\tau+b\) have actual envelopes \(F(P_A(p^*b))\), continuous
relative to the reference on the divisor complement, and with the
prescribed all-point poles.

\subsection{9.2. Both derivative signs in every original
direction}\label{both-derivative-signs-in-every-original-direction}

Applying G0 to \(w_t=\psi+t p^*f\) and (9.2)--(9.4) gives
\[\left.\frac d{dt}\right|_{t=0}\mathscr O_{\mathcal M}(u+tf)
=\int_Y p^*f\,M_A(\psi)=\int_X f\,\widetilde M_{\mathcal M}(u). \tag{9.5}\]
This is a two-sided derivative for every \(f\in C^0(X^{\rm an})\), not
just directions descending from a projective contraction or PL tests.
Its factor is raw mass \(V\). For clarity, when extending a PL
derivative by uniform approximation with error \(\epsilon\), the
obstacle bounds are \(\pm|t|\epsilon\) and the quotient errors are
\(\pm V\epsilon\). Replacing \(f\) by \(-f\) gives the second sign; no
\(t^2\epsilon\) bound is used.

\subsection{9.3. An actual attained maximizer and the conditional
equation}\label{an-actual-attained-maximizer-and-the-conditional-equation}

Let \(\mu\) be a positive Radon measure of mass \(V\) satisfying (7.1).
Its unique Radon inverse \(\nu\) has finite ample dual energy by (7.2).
G0 gives an actually attained genuine \(\psi\in\mathcal E^1(Y,A)\)
maximizing \(E_A(\eta)-\int\eta d\nu\). Hence \(u=F(\psi)\) is an actual
original metric attaining (7.1), not a completion or an approximate
object.

The mechanism can be expressed on the original carrier. A normalized
maximizing family has bounded ample \(J\), hence bounded original \(J\)
under (5.4). Original normalized compactness gives a subnet to an actual
\(u\); energy usc and the ample bounded-\(J\) target \(L^1\) continuity
through (7.2) give attainment. Nonmetrizable compactness is used as
subnet compactness. No countability or countable field containing an
entire psh family is required.

For every original continuous \(f\) and real \(t\), the genuine minorant
\(v_t=F(P_A(w_t))\) satisfies \(q_{v_t}\leq q_u+tf\), by order and
(9.1). Integrability and the maximal objective therefore give
\[\mathscr O_{\mathcal M}(u+tf)-\int(q_u+tf)d\mu
\leq E_{\mathcal M}(u)-\int q_ud\mu.\] Use (9.5), first with \(t>0\) and
then with \(t<0\), reversing the inequality when dividing by a negative
number. We obtain equality of the two finite Radon measures on every
continuous original test:
\[\widetilde M_{\mathcal M}(u)=\mu.                       \tag{9.6}\]
This is an actual equation for the pairing (8.1). It is not yet labeled
the genuine original local Monge--Ampere equation because Section 10
remains.

For two solutions of (9.6), their ample measures have zero mass on \(D\)
and identical pushforwards. Radon inversion proves equality of those
ample measures; G0 uniqueness makes their genuine quotients differ by
one constant. Thus their actual original metrics differ by the same
constant at every point, including their prescribed poles. The converse,
that such a solution has finite dual energy and computes the dual
supremum, follows from the ample energy concavity statement through
(7.2).

\subsection{9.4. Comparison and domination for the
pairing}\label{comparison-and-domination-for-the-pairing}

For \(u,v\in\mathcal E^1_{\mathcal M}\), the strict sets outside \(P\)
correspond exactly to the strict quotient sets outside \(D\), and both
product measures charge \(P\) by zero. Ample comparison therefore gives
\[\int_{\{u<v\}}\widetilde M_{\mathcal M}(v)
\leq\int_{\{u<v\}}\widetilde M_{\mathcal M}(u).            \tag{9.7}\]
Strict-winning locality and mixed locality transport in the same way. If
a genuine \(u\in\mathcal P_{\mathcal M}\) is at most a finite-energy
\(v\) almost everywhere for \(\widetilde M(v)\), Radon inversion
transfers that inequality off \(D\) to the ample measure. It gives \(D\)
zero mass, so ample domination applies and yields \(Q(u)\leq Q(v)\)
everywhere. Then \(u\leq v\) everywhere. None of these deductions
assumes projectivity of the original carrier or a prescribed skeletal
support.

\section{10. The exact unclosed local-operator
statement}\label{the-exact-unclosed-local-operator-statement}

Sections 8.2--8.4 prove the following restricted local statement: the
actual local finite-section products converge for divided global-section
nets on compact original frame neighborhoods off \(P\); their bounded
endpoint limits are germ independent inside that category; fixed
analytic-unit frame changes pass through the entire decreasing
construction; relative finite-energy floor restrictions and integrable
weighted limits recover the raw pairing (8.1).

There is a useful exact frame reinterpretation. On \(Y\setminus D\),
division by the nonvanishing divisor section gives the actual local
bundle isomorphism
\[\pi^*L\longrightarrow A,\qquad e_L\longmapsto e_A/f.\] The original
raw weight in \(e_L\) is therefore literally the raw weight of the
genuine \(A\) metric in the actual frame \(e_A/f\). This is an isometry,
not a numerical replacement of a reference. It may permit a more direct
comparison with the actual-ample construction and avoid globalization of
an arbitrary local contact family. The present proof has verified that
comparison only through the specific divided global-section limits in
Section 8; it has not audited the full bounded and unbounded local
construction under this identification.

What is not proved is the following \textbf{local bridge L}.

\begin{quote}
Let \(Q\) be a compact original frame neighborhood, with ordinary
interior, whose closure lies in \(X^{\rm an}\setminus P^{\rm an}\). For
a bounded genuine \(A\)-psh tuple, allow the local rational contact
approximants and actual ample collars used in the original local
intersection construction. For the literal raw weights
\(\lambda_i=\psi_i+h_A+\log|f|\) on \(Q\), prove that this local
construction is independent of its contact approximants, collar
realizations, and compatible frames, and equals the divided
global-section limit in (8.4). Establish this in its bounded, possibly
discontinuous category, with its integrable weighted version. Then prove
that its finite-energy, relative-floor nonpluripolar extension agrees
with (8.6).
\end{quote}

This statement concerns only the actual metrics and approximation
category needed here. It does not ask for arbitrary local approximations
without common neighborhoods, arbitrary weak convergence of psh
functions, or choice independence for different retained
\((L,\pi,A,D)\).

The precise logical issue is that G0's global bounded decreasing
convergence identifies the limit \textbf{after} the local approximants
have been realized as globally genuine potentials in the fixed ample
class. Section 8 constructs one sufficient original-section family and
proves independence among such globally realizable families. A merely
local contact family supplied to a collar need not itself be a
decreasing family of global fixed-\(A\) potentials. One route would
prove simultaneous reference-preserving globalization and compare the
resulting products. Another would verify that the exact actual-ample
local construction in the frame \(e_A/f\) already supplies this
comparison, using the preceding isometry. This is a scope check for the
actual constructed metrics, not a demand for a larger new local psh
theory. Neither route has been completed here. BE's uniform theorem for
continuous endpoints does not fill it.

The pertinent newly opened primary proof is Wu arXiv:2609.27610v1,
Proposition 3.28 and Theorem 4.2. Its printed collar argument uses local
PL positivity, a global PL-to-psh conversion, and a common collar before
monotone product convergence. Its printed formulas also contain
reference slips. With a collar function \(v=h_0+p\) and scale \(a>0\),
if \(\Psi\) is the glued scalar in the metric reference of \(v\), then
the actual canonical-reference potential is \[\Phi=p+\Psi/a,\] and on
the inner open
\[\lambda=a\Phi+a h_0+M+1.                              \tag{10.1}\] The
\(p\) term cannot be omitted or treated as a numerical bundle change.
Furthermore, a maximum containing a zero branch cannot have a negative
value at the collar's tested point. These literal reference points must
be repaired in any application of that proof. The source's local
definitions and proofs are printed for sequences; an application over
arbitrary-cardinality \(k\) must retain nets or separately justify every
sequence replacement. The supplied ample G0 conclusions do not grant
these new local globalization steps by their historical labels.

The continuous finite-stage identities, the compact reference bounds,
and the ample mixed weighted limits needed for a bridge are available. A
plausible approach is to construct the finite collar weights with the
literal identity (10.1), prove their actual nef/model/FS membership in
one fixed ample bundle, and apply finite inner approximant flattening
\textbf{before} adding the common polar section. This paper has not
completed that argument. We consequently retain L as a missing
mathematical lemma. No counterexample to L or to P16.2 is established.

If L is proved, Sections 3--9 supply the rest of the retained
presentation theory in the full stated domain: the original section
cone, energy, mass \(A^n\), intrinsic integrability criterion, every
continuous original derivative, an attained solution, comparison, and
uniqueness modulo constants. Without L, (9.6) remains a pairing equation
with extensive local evidence, and is not promoted to the exact required
original-carrier equation theorem.

\section{11. A genuinely nonprojective test over every retained
field}\label{a-genuinely-nonprojective-test-over-every-retained-field}

The geometric domain is not empty. The following is the classical
Miyake--Oda example, printed in Fujino--Sato Example 4.4, not a new
example and not a restriction of the domain of Sections 3--9.

In \(N=\mathbb Z^3\) use \[
v_1=e_1,\ v_2=e_2,\ v_3=e_3,\ v_4=(-1,-1,-1),\quad
v_5=(0,-1,-1),\ v_6=(-1,0,-1),\ v_7=(-1,-1,0).
\] The maximal cones are
\[123,\ 125,\ 256,\ 456,\ 236,\ 367,\ 467,\ 137,\ 157,\ 457.\] Start
with the four maximal cones of the fan of \(\mathbb P^3\). For an
ordered old basis \((a,b,c)\), the two additional rays are \(a+c,b+c\).
The three cones with bases \((a,b,a+c)\), \((b,a+c,b+c)\),
\((c,a+c,b+c)\) cover the nonnegative old octant in the regions
\(x\geq z\), \(x\leq z\leq x+y\), and \(z\geq x+y\), respectively. Their
shared faces match. Apply this to the ordered bases \((v_1,v_2,v_4)\),
\((v_2,v_3,v_4)\), and \((v_3,v_1,v_4)\). This gives the ten listed
cones and a complete fan. The new determinants are \(1,-1,1\) relative
to the old unimodular basis. Every maximal chart is therefore
\(\mathbb A^3_k\), over any field.

For completeness, the primary properness theorem we use is Rohrer,
dissertation IV 1.3.11--1.3.12, whose printed proof extends a DVR-valued
character into a cone of the appropriate complete star fan. Star
completeness itself has a direct proof that avoids a misleading
origin-based argument: choose \(z\) in the relative interior of the
fixed cone \(\sigma\); along a sequence \(z+\epsilon_j x\) with
\(\epsilon_j\downarrow0\), completeness and finiteness give one cone
\(\tau\) containing infinitely many terms and \(z\). Its intersection
with \(\sigma\) contains an interior point, so \(\sigma\) is a face of
\(\tau\). Projecting gives \(x\) in \(\tau/\sigma\). Thus the star fan
is complete. The DVR extension proof then establishes properness of the
toric carrier. Separatedness follows from the fan gluing theorem in the
same primary source. The charts and their shared dense torus prove
smoothness and integrality.

Every line bundle can be represented by an invariant boundary Cartier
divisor \(\sum a_iD_i\): on the dense torus the Laurent polynomial ring
is a UFD, so the line bundle is trivial there; take its rational
trivializing section and its boundary divisor. The three invariant wall
curves for \(25\), \(36\), and \(17\) are \(\mathbb P^1\), since their
quotient fans have the two primitive opposite rays and the corresponding
charts glue by inversion. Their degrees are \[
d_{25}=a_1+a_6-a_2-a_5,\qquad
d_{36}=a_2+a_7-a_3-a_6,\qquad
d_{17}=a_3+a_5-a_1-a_7.                                \tag{11.1}
\] For example, the wall relation is \(v_1+v_6=v_2+v_5\); the local
Cartier frames have characters pairing to \(-a_i\) with the respective
rays, and their difference has transition exponent \(a_1+a_6-a_2-a_5\)
on the quotient coordinate. The rational first frame section has that
order at the second chart's origin, giving the stated sign. The other
two relations are cyclic. The three numbers in (11.1) sum to zero. An
ample line bundle would have positive degree on all three proper curves,
which is impossible. Proper plus ample implies projective (Stacks 0B5N),
so this smooth proper carrier is genuinely nonprojective over every
field in our domain.

Section 3 applies to this example, as to every smooth integral proper
\(X\) and every chosen big actual rational \(L\). No toric invariance or
support restriction is used in the metric or pairing construction.

\section{12. External-premise ledger and literal
evidence}\label{external-premise-ledger-and-literal-evidence}

The machine-readable register \texttt{external-premises.json} is part of
this paper. Each record states the precise mathematical premise, the
carrier and field scope, the exact local source unit, its original URL
and supplied coordinates when present, the local literal coordinates and
hash, and whether its whole work or recursively cited proofs were read.
The latter flags are false unless explicitly stated. Original historical
filesystem paths are metadata only; no such paths were opened.

The essential premises and their application boundaries are as follows.

\begin{enumerate}
\def\labelenumi{\arabic{enumi}.}
\tightlist
\item
  \textbf{G0.} The exact supplied actual-ample theorem is used on smooth
  projective \(Y\) and actual ample rational bundles only. Its source is
  the preserved supplied body \texttt{evidence/packet-00.txt}. No
  nonample theorem from its underlying citations is granted separately.
\item
  \textbf{Chow and resolution.} Stacks 0200/0201 supply the proper
  birational projective refinement; Temkin v2 Theorem 1.2.1,
  finite-sequence definitions, stabilization lemma, and pertinent main
  proof supply a finite regular-final resolution. Stacks 07QW and 0B8X
  verify excellence and perfect-field smoothness. The omitted
  localization proof in 07QU is not used. Earlier resolution algorithms
  cited by Temkin have not all been read.
\item
  \textbf{Actual sections and fibers.} Stacks 0AY8 and 0BFP are applied
  to normal original \(X\), proper birational \(Y\to X\), reduced
  integral \(Y\), and the common generic function field. Their pertinent
  printed statement/proof units are preserved. Stacks 0AG0 supplies
  regular local factoriality. Stacks 01Q3, 0B5T, and 0BEM supply global
  generation and coherent growth, with their full pertinent printed
  proofs retained. Associated-point regularity is proved by the
  nonzerodivisor criterion in the actual finite module; Stacks 0B3Q's
  printed omitted proof is not described as read.
\item
  \textbf{Analytification.} Berkovich 1993, 2.6.1/2.6.7/2.6.8 give
  functoriality, products/open compatibility, separatedness, and the
  image/valuation extension statement. Compactness, closed quotient
  descent, center compatibility, and the canonical original metric
  application are reconstructed in Section 4. We do not invoke 2.6.9's
  cross-referenced properness proof.
\item
  \textbf{FS arithmetic and fixed-divisor foundations.} BJ v4
  Definitions 2.18/4.1 and section arithmetic 2.19--2.23 are used
  literally. Ample determination, compactness, and divisorial envelope
  principles are used only in their actual ample G0 scope, with the
  supplied BJ proof units retained for source checking. Nonample
  rational approximation, homogeneous projection, decomposition,
  division, and the zero-class argument are proved in Section 5. BJ 4.6
  and dependent nonample applications are diagnostic evidence, not
  premises.
\item
  \textbf{Actual continuous products.} BE v3 7.14 and 8.3--8.6, 8.18,
  and BJ Remark 7.23 supply continuous actual-metric compatibility and
  raw fundamental-cycle products. CLD v2 18.1.2 and 8.8.5 supply literal
  frame and zero-Hessian identities at continuous stages. CLD 8.7.2 is a
  uniform-convergence statement; it is not a bounded discontinuous or
  unbounded local bridge. The noncontinuous parts proved here use G0
  decreasing and weighted convergence; the missing category comparison
  is labeled L.
\item
  \textbf{Radon topology.} Compact Hausdorff normality, Tychonoff
  compactness, Riesz representation, regularity, and extension of a
  complete-field norm to a finite extension are standard background,
  also used in the published ample theory. Their new applications and
  measure uniqueness are proved here. No metrizable topology or
  countable selection is invoked.
\item
  \textbf{Toric test.} Fujino--Sato Example 4.4 records the example. The
  precise complete-fan properness proof is from Rohrer's dissertation,
  not the separate paper that only cross-refers to it. The chart, star,
  and degree applications are reconstructed in Section 11. These
  premises are unnecessary for the whole-domain construction, which uses
  Section 3 instead.
\end{enumerate}

The newly retrieved public renderer evidence is preserved in the
directory \texttt{evidence/}. Exact response JSON and text files use the
prefix \texttt{public-}. The attempted shell download of the Wu PDF
failed with exit code 7; its stderr is retained. No PDF bytes or
original-file checksum were invented. Local literal LF coordinates are
distinct from original renderer labels. All 18 supplied packet blocks
and 124 individual supplied literal units are retained without
consulting any historical worker filesystem.

\section{13. Reproduction and exact negative
conclusion}\label{reproduction-and-exact-negative-conclusion}

Run \texttt{python3\ rerun.py} in this workspace. It re-extracts the
immutable input, preserves and verifies literal supplied evidence,
checks the exact status schema and backwards note supports, records
Python/Pandoc/XeLaTeX versions and a minimal platform description, runs
the small rational diagnostic calculations, and builds
\texttt{P16.2-paper.tex} and \texttt{P16.2-paper.pdf} using the existing
tools. Exact argument arrays and stdout/stderr are retained in
\texttt{logs/}. The generated \texttt{artifact-manifest.json} hashes the
local artifacts. No packages are installed, and no heavy computation is
launched.

The mathematical outcome is the proof reconstruction in Sections 2--9
and 11, plus the precise unproved bridge L in Section 10. In particular,
the unsafe approximation theorem has not been restored by a historical
label, and the pairing equation has not been relabeled as the required
genuine original-carrier Monge--Ampere equation. The standalone paper,
its literal ledger, and the conditional deductions are the
implementation deliverable. Original P16.2 remains unproved by this
work.

\end{document}
